% fig-d5-minor-ladder.tex --- D5, the determinantal ladder for a Seidel matrix.
%
% Defines \ClebschMinorLadder.
%
% Content is the paragraph "The same statement in principal minors" in
% papers/clebsch-passages/sections/05-golden-operator.tex, together with
% thm:aligned-faithfulness.  The point of the picture is that the
% reconstruction theorem deliberately consults the weakest rung.

\providecommand{\ClebschMinorLadder}{%
\begin{tikzpicture}[every node/.style={outer sep=0pt}]

\node[csfnode,align=left,text width=84mm] (m1) at (0,0)
  {\textbf{order 1.}\;\(\det C[\{i\}]=0\).  Every first-order principal
   minor vanishes: the diagonal is zero.  No information.};

\node[csfnode,align=left,text width=84mm] (m2) at (0,-1.35)
  {\textbf{order 2.}\;\(\det C[\{i,j\}]=-1\), constantly.  No
   information.};

\node[csfnode,align=left,text width=84mm] (m3) at (0,-3.0)
  {\textbf{order 3.}\;\(\det C[\{i,j,k\}]=2\,\epsilon(ijk)\), twice the
   triangle sign.  The full third-order list hands over every triangle
   product, so reconstruction up to switching is immediate: this rung
   trivializes the problem.};

\node[csfnode,align=left,line width=1pt,fill=black!8,text width=84mm] (m4) at (0,-5.35)
  {\textbf{order 4.}\;\(\det C[Q]=3-2w(Q)\in\{-3,5\}\): one bit per
   four-set, and invariant under \(C\mapsto-C\).  This is the data the
   reconstruction theorem uses.  It determines \(C\) up to switching and
   global negation, it needs at least seven vertices --- and seven is
   sharp --- and \(O(n^{2})\) selected determinants suffice.};

\draw[decorate,decoration={brace,amplitude=5pt},semithick]
  ($(m1.north east)+(2mm,0)$) -- ($(m3.south east)+(2mm,0)$);
\node[anchor=west,align=left,font=\scriptsize,text width=36mm]
  at (5.0,-1.5)
  {strictly stronger data.\\ The theorem consults none of it.};

\draw[csfarr] (5.0,-4.6) -- ($(m4.east)+(2mm,0.4)$);
\node[anchor=west,align=left,font=\scriptsize,text width=36mm]
  at (5.0,-5.6)
  {the weakest rung, and the one the theorem uses};

\end{tikzpicture}%
}
